\documentclass[11pt]{article}
\usepackage[a4paper,margin=25mm]{geometry}
\usepackage{amsmath,amssymb,amsthm,booktabs,xurl,hyperref}
\hypersetup{hidelinks}
\setlength{\emergencystretch}{3em}
\hyphenation{two}
\newtheorem{theorem}{Theorem}
\newtheorem{lemma}[theorem]{Lemma}
\newcommand{\Q}{\mathbb Q}\newcommand{\Z}{\mathbb Z}
\title{Removing the c-integer predicate from the combined quartics}
\author{Z diophantine in Q project}
\date{24 November 2026}
\begin{document}
\maketitle

\begin{abstract}
The proved half-trace integrality theorem replaces one complete
integer predicate in each of the current general and specified
quadratic systems. After preserving the outer source's square
relations and updating the actual shared contact sum, the complete
systems have 887 variables and 750 quadratic equations, or 1052
variables and 913 equations. These improve the 943/798 and
1124/975 predecessors by 56 and 72 variables. The general system
restores one half-trace square; the specified system already
retains it. All remaining poles keep their BB-native interpretation.
The equivalence concerns the coded-instance projection under the
retained arithmetic interfaces. It is not a bijection of complete
auxiliary zero sets.
\end{abstract}

\section{The arithmetic replacement and its order of use}

We use the proved theorem in \cite{replacement}, including its
signed rational Pell classification. Its exact statement is:
\begin{theorem}[Half-trace integrality]\label{arithmetic}
Suppose $a\in3\Z$, $c,d,\zeta\in\Q$, and
\[
 d^2-(a^2-1)c^2=1,\qquad
 T=d-(a-2)c\in\Z,\qquad
 \zeta^2=2\bigl(ad+(a^2-1)c+1\bigr).
\]
Then $c,d\in\Z$. The assertion includes $a=0$.
\end{theorem}

For completeness, the denominator mechanism is as follows. Set
$D=4a-5$ and $L_a=2a^2-2a-1$. Substitution of
$d=(a-2)c+T$ gives
\[
 (Dc-(a-2)T)^2=(a^2-1)T^2-D.
\]
Its right side is integral. A rational square root of an integer
is integral, so $Dc\in\Z$. If $N$ is the positive reduced
denominator of $c$, then $N\mid|D|$. Since
\[
                 8L_a+3=D(4a+1)
\]
and $3\mid a$, one has $\gcd(2L_a,N)=1$. The half-trace square
is $2L_a c+2aT+2$, whose reduced denominator is therefore
exactly $N$. It follows that $N$ is a square.

The rational Pell classification in \cite{replacement} now excludes
every nonintegral orbit. For $a\ge3$ divisible by three, a
nonintegral orbit has exact denominator $4a-5\equiv3\pmod4$.
For $a=-A\le-3$ with $3\mid A$, it has exact denominator
$4A+5\equiv2\pmod3$; the additional seed orbit requires $4A+5$
to be a square and is unavailable. At $a=0$, the only
nonintegral denominator is five. Thus none has square denominator.
This uses the proved classification, not a finite search over Pell
solutions.

\subsection{The source gives the hypotheses before c is certified}

The retained preliminary source reconstruction yields integral
$q,S_0,T_0,w,s,b,\gamma$ without using $c\in\Z$.
For the general route this is the preliminary auxiliary-power and
prefix argument of \cite[Section~2]{signed}; for the specified
route $q$ is given by the retained explicit integer power circuit.
The literal source equations give
\[
 n=q^{16},\quad M=n^2=q^{32},\quad
 r=S_0(M-n)+T_0(M-1),\quad
 a=(wM+1)rsM.
\]
Modulo three, either $M\equiv0$, or $M\equiv n\equiv1$.
In the latter case $r\equiv0$. Hence $3\mid Mr$ and $3\mid a$,
before a final identity for the power or a first-Pell index is used.

The retained linear row is
\[
                    d=(a-2)c+bw+(4a-5)\gamma,
\]
so its shifted trace is integral. The source already stores
$A_c=ac$. The quadratic half-trace row is
\begin{equation}\label{half}
             \zeta^2-2ad-2aA_c+2c-2=0.
\end{equation}
It is named \texttt{jones1980.pell\_positive\_root}, and its
coordinate is \texttt{pell\_d\_positive\_square}.
The general predecessor had removed this row in v101.
The specified predecessor still contains both the row and coordinate.
Thus only the general system needs to add them back.

Theorem~\ref{arithmetic} now proves $c,d\in\Z$ before any later
use of integral first-Pell classification. In particular the
zero case $a=0$ is handled before the source excludes $s=0$.
This avoids assuming the very integrality supplied by the deleted
predicate.

\subsection{Canonical completeness}

For a positive canonical coding witness the retained source has
$a\ge2$, $B=2r+1$, $c=Y_B(a)$ and $d=X_B(a)$, with the
standard Pell sequences. Their identities give
\[
 ad+(a^2-1)c=X_{B+1}(a)
            =X_{2r+2}(a)=2X_{r+1}(a)^2-1.
\]
Consequently $\zeta=2X_{r+1}(a)$ supplies \eqref{half}.
The existing positive source witness can therefore be chosen
before constructing the smaller tuple of integer-interface
witnesses. This preserves the coded instance, not every
arbitrary signed auxiliary witness of the predecessor.

\section{Removing the borrowed square without losing a source equation}

The c-integer block supplied a helper $X_c$, literally
\texttt{int\_c\_37}, with
\begin{equation}\label{square}
                             c^2-\rho X_c=0.
\end{equation}
Here $\rho$ is the common contact gauge and the retained
$\rho\epsilon=1$ guard makes it nonzero. Once the c-interface
rows are removed, the only remaining occurrences of $X_c$ are
\[
 d^2-A_c^2+\rho X_c-1=0,\qquad Ev-\rho X_c=0,
\]
where $v=\texttt{i\_inverse}$ belongs to the retained reciprocal
integer input. Replace these rows by
\begin{equation}\label{outer}
                 d^2-A_c^2+c^2-1=0,\qquad Ev-c^2=0.
\end{equation}
They remain quadratic. The first new residual is the old
residual plus \eqref{square}; the second is the old residual
minus \eqref{square}. Conversely $X_c=\epsilon c^2$ recovers
all three old equations from \eqref{outer} and the gauge guard.
This exact elimination divides by neither $c$ nor $v$ and
preserves the case $c=0$.

The definition $A_c-ac=0$ is retained literally. Thus the first
equation of \eqref{outer} is exactly the first Pell norm used by
Theorem~\ref{arithmetic}. We remove the square helper and its
defining row, saving one more coordinate and row than a version
that retains this borrowed source helper.

\section{The actual contact sum after deleting one predicate}

All coordinates whose names start with \texttt{int\_c\_} are
removed. The c-private rows have prefixes
\[
 \texttt{integer.c.},\quad \texttt{global.c.},\quad
 \texttt{scalar\_norm.c.},\quad \texttt{detector\_trace.c.}.
\]
These include its point, scalar, contact, comparison, and pole
equations. Common reference, numerator, scalar, detector, and
contact-coefficient rows remain. In particular the root b-block
and the common spacing $D=T_b/16$ remain unchanged.

The compiled predecessors have already eliminated the final
contact weight of the reciprocal block using the shared sum.
In the stored normalization, if $\omega_*$ denotes this omitted
last weight, the identity is
\begin{equation}\label{oldsum}
                   \omega_*=\epsilon-\sum_{\nu\ne *}\omega_\nu.
\end{equation}
The common physical rescaling of all these weights is unchanged.
The c-weights removed from the sum are
\begin{center}
\begin{tabular}{@{}ll@{}}
\toprule
General & \texttt{int\_c\_153}, \texttt{int\_c\_last\_contact\_weight}\\
Specified & \texttt{int\_c\_154}, \texttt{int\_c\_155},
 \texttt{int\_c\_last\_contact\_weight}\\
\bottomrule
\end{tabular}
\end{center}
The new expression is
\begin{equation}\label{newsum}
 \omega_*^{\rm new}
   =\epsilon-\sum_{\substack{\nu\ne *\\\nu\notin c}}\omega_\nu
   =\omega_*^{\rm old}+\sum_{\nu\in c}\omega_\nu.
\end{equation}
This is a new shared-gauge condition for the reduced tuple.
It does not assert that an omitted physical contact function
vanishes, nor supply an extension to the old full gauge sum.

There are exactly three retained occurrences of the eliminated
weight. Write $\epsilon_2$ for the existing square register,
$B_*,G_*,V_*,D_*$ for the final reciprocal pole input, gate,
inverse-weight register, and norm error. The three rows before
substitution are
\begin{align}
 167281B_*(\epsilon_2+\omega_*)-166464G_*^2&=0,\label{pole}\\
 \epsilon\epsilon_2-\omega_*V_*&=0,\label{inverse}
\end{align}
and, in the general route,
\begin{equation}\label{ggate}
 2\epsilon\epsilon_2-\epsilon V_*-\epsilon_2^2
          +2\epsilon_2G_*-\epsilon_2\omega_*-D_*^2=0,
\end{equation}
or, in the specified route,
\begin{equation}\label{sgate}
 2\epsilon G_*+2\epsilon-\epsilon_2-\rho V_*-\omega_*-D_*^2=0.
\end{equation}
Their local index dictionaries are
\begin{center}
\begin{tabular}{@{}lrrrr@{}}
\toprule
Interface & $B_*$ & $G_*$ & $V_*$ & $D_*$\\
\midrule
General &73&72&158&156\\
Specified &174&173&161&158\\
\bottomrule
\end{tabular}
\end{center}
Every index uses the \texttt{int\_i\_reciprocal} prefix.
The checker derives both old and new rows from these templates
and \eqref{oldsum}--\eqref{newsum}. Equivalently, for compilation
one deletes the c-weight monomials in those three expanded rows.
No other retained row contains a deleted contact weight.

\subsection{Why the same interface theorem applies}

The global-gauge proof \cite{global,bb} works for any fixed finite
tuple of integer inputs. Its soundness lemma uses the identity
$\rho=q_j+(\rho-q_j)$ for each remaining contact, and its
completeness construction applies the finite-family prime supplier
to the remaining distinct contacts. Deleting the c-contact lists
introduces no new collision or denominator condition. The retained
first-contact indices are the previous index list with index seven
removed; the labels of the other contacts are not renumbered.
The sixteenth endpoint and every interpolation definition remain.

The physical sum now has 24 comparison weights in the general
route and 30 in the specified route. At a selected prime only
the selected weight has its prescribed pole, so the new gauge
still has valuation $-1$. Foreign-contact parity and the height
estimates are unchanged. At the fixed prepared primes and at
real infinity the same finite-family asymptotic is
\[
 q_j=\kappa\lambda^{-2}+O(\lambda^{-3}),\qquad
 \rho=m\kappa\lambda^{-2}+O(\lambda^{-3}),\quad m=24\text{ or }30,
\]
and
\[
 \frac{(\rho-q_j)^2}{q_j}
       =(m-1)^2\kappa\lambda^{-2}+O(\lambda^{-3})\longrightarrow0.
\]
The proof divides by neither $m$ nor $m-1$. The same common
finite neighborhoods and shared scalar-depth construction therefore
give BB-prepared comparison and scalar inputs for the reduced
tuple. Its 36 or 40 private pole blocks retain the coefficients
and interpretation from \cite{bb}.

For completeness all these interface witnesses are chosen anew
after the canonical coding data are fixed. For soundness the
remaining predicates first certify their own distinguished inputs;
the source argument in Section~1 then supplies integrality of $c$.
There is no attempted projection of a new auxiliary zero into an
old c-certificate with the former gauge.

\section{Complete systems, counts, and independent verification}

\begin{theorem}\label{main}
Under the same retained non-pole arithmetic interfaces, BB-native
pole theorem, and external Jones coding theorem as the preferred
v105 systems, the new complete systems have the same coded-instance
projection. Their sums of squares are quartics in 887 and 1052
rational variables, respectively. No new arithmetic hypothesis or
numerical height threshold is introduced.
\end{theorem}
\begin{proof}
Theorem~\ref{arithmetic}, \eqref{outer}, and the remaining
integer predicates give source soundness. Canonical completion
of \eqref{half} and the reduced simultaneous interface construction
give source-instance completeness. Every listed residual is an
ordinary quadratic; its sum of squares vanishes over $\Q$ exactly
when every residual vanishes. Nonzero quadratic leading forms
make the degree exactly four.
\end{proof}

The complete census is
\begin{center}
\begin{tabular}{@{}lrrrrrr@{}}
\toprule
Interface & Old $N$ & Removed & Added & New $N$ & Old rows & New rows\\
\midrule
General &943&57&1&887&798&750\\
Specified &1124&72&0&1052&975&913\\
\bottomrule
\end{tabular}
\end{center}
The general system deletes 49 rows, changes five, and adds the
one half-trace row. Its other 744 rows are identical up to
coordinate renumbering. The specified system deletes 62 rows,
changes five, and has 908 unchanged rows, including its existing
half-trace row. The resulting predicate counts are twelve and
ten, respectively. They include the retained reciprocal input
in each route.

The independent verifier imports no compiler code. It pins both
predecessor SHA-256 hashes, audits every occurrence of all 57 or
72 removed coordinates, checks the two source-square identities,
the inverse-gauge guards, the source linear and product rows,
and the exact half-trace polynomial. It derives all three changed
contact rows independently from the omitted-weight formulas.
Every other row is compared by named monomials. It also checks
the updated predicate, contact and pole metadata, complete variable
support, coefficient rings, degree bounds, and output hashes.

The separate arithmetic theorem has its own proof and verifier
\cite{replacement}. These finite checks do not prove the retained
height, prime supplier, elliptic carrier, BB, or coding theorems.
Earlier hashes and coordinate maps in inherited metadata describe
historical stages. The current certificate explicitly records
source-instance equivalence and rejects an old/new zero-set
bijection claim.

\begin{thebibliography}{9}
\bibitem{replacement} \emph{A half-trace square replaces the
c-integrality predicate}, including the signed rational Pell
classification and its separate audit.
\nolinkurl{undecidability/pell_integrality_replacement_v106/pell_integrality_replacement.tex}.
\bibitem{signed} \emph{Both first-Pell signs have the same coding
projection}.
\nolinkurl{undecidability/pell_signed_branch_v101/signed_branch_projection.tex}.
\bibitem{input} \emph{Sharing squared input comparisons in the
complete quartic systems}.
\nolinkurl{undecidability/input_comparison_v105/input_comparison_sharing.tex}.
\bibitem{global} \emph{One contact gauge in a combined quadratic system}.
\nolinkurl{undecidability/papers/global_gauge_quartic_checkpoint.tex}.
\bibitem{bb} \emph{Congruent-number poles in the complete combined
quartic systems}.
\nolinkurl{undecidability/quadratic_reduction_v104/bb_native_combined.tex}.
\end{thebibliography}
\end{document}

